% DRAFT: fresh native jets and finite-sieve completion remain pending.
% Input work/deg10/next_cases_20261030/dyadic_formal_kernel.tex beforehand.
\subsection{Excluding a point of order $51$}
\begin{proposition}\label{prop:decic-no51}
An elliptic curve over a number field of degree $10$ has no point of order $51$.
\end{proposition}
\begin{proof}
Put $X=X_1(51)$ and $J=J_1(51)$. The curve has genus $65$, good
reduction at $2$, and $\operatorname{rank}J(\mathbb Q)=0$ by
\cite[Theorem~3.1(2)]{DEHMZ}.
A proposed torsion point gives a rational effective noncuspidal
divisor $D$ of degree $10$ on $X$, taking its trace with
multiplicity if its moduli point is defined over a proper subfield.

Independent exhaustive enumeration of all elliptic curves,
twists and level structures over $\mathbb F_{2^f}$, $f\leq10$,
shows that the only noncuspidal closed-point degrees in this
range are $9$ and $10$. There are respectively $80$ and $64$
points, forming five and four diamond orbits.
For every representative $P$, exact Riemann--Roch on the full
curve gives that
\[
 (T_5-1-5\langle5\rangle)P
\]
is nonprincipal. The six untouched native $5$-isogenies at
each representative independently reproduce the finite
correspondence, including its closed-point degrees and norm
multiplicities.

The operator $A_5=T_5-1-5\langle5\rangle$ kills $J(\mathbb Q)$,
by rank zero, Eichler--Shimura at $5$, and injectivity of rational
torsion reduction at $5$. It kills each rational cusp as a
literal divisor, by the Tate parametrization and diamond
transitivity. A degree-$10$ reduction with noncuspidal support
is one degree-$10$ point or a degree-$9$ point plus a rational
cusp. Either possibility contradicts the computed
nonprincipality. Thus $\overline D$ is entirely cuspidal.

The cusp signature at $2$ is $1^{16},2^8,8^4$; the $j$-map has
degree $1152$. Put
\[
 U=(\mathcal C-1)H^0(X_{\mathbb F_2},\Omega^1),\qquad\dim U=62.
\]
We verify that restriction from $U$ has rank $10$ at every
cuspidal effective divisor of degree $10$.
The existing common-basis native completion data provide jets
through multiplicities $8$, $4$, and $1$ at the rational,
quadratic, and octic cusps respectively. These restrictions are
independent of the total degree of the divisor.
An independent binary replay checks all $9655880$ degree-$10$
divisors satisfying these individual bounds, and every rank is $10$.
The omitted divisors are exactly
\[
 10P,\qquad9P+Q\ (P\ne Q),\qquad5R,
\]
with $P,Q$ rational cusps and $R$ quadratic. There are
$16+16\cdot15+8=264$ of them.

The diamond $\langle5\rangle$ is transitive on the $16$ rational
cusps and on the eight quadratic cusps, as checked on primitive
Tate cusp labels. Thus it suffices to fix $P$ and check $10P$,
the fifteen divisors $9P+Q$, and one divisor $5R$.
Fresh exact native completions give the restrictions through
order ten at $P$ and order five at $R$, with enough precision
for every coefficient used. The ranks for $10P$ and $5R$ are $10$.

To combine the new ninth-order jet with the recorded first-order
jets at the other rational cusps, we identify the differential
bases explicitly. Let $C_{\rm old},C_{\rm new}$ be the two Cartier
matrices, and let $W_{\rm old},W_{\rm new}$ be the kernels of
restriction at $8P$. The equations
\[
 C_{\rm old}T=TC_{\rm new},\qquad
 W_{\rm old}T\subseteq W_{\rm new}
\]
have a one-dimensional solution space over $\mathbb F_2$.
Its nonzero matrix $T$ is invertible and takes the old Cartier
image to the new one. A geometric solution exists after a
diamond identifies the two choices of rational base cusp.
Uniqueness therefore proves that the computed transition is
correct up to that diamond. It transports the complete set of
rational first-order restrictions, so all fifteen mixed ranks
are checked in a single basis. Each is $10$.
An independent binary calculation checks the dimension-one
intertwiner space, its invertibility, the anchor and Cartier-image
identities, every transported restriction, and all seventeen
new ranks. Hence all
\[
 [t^{10}](1-t)^{-16}(1-t^2)^{-8}(1-t^8)^{-4}=9656144
\]
cuspidal divisors satisfy the formal-immersion condition.

Apply Corollary~\ref{cor:dyadic-Cartier-formal-immersion} to the
endomorphism $T_2-\langle2\rangle$ of the full Jacobian.
The etale cuspidal lift $C$ of $\overline D$ is a
$\mathbb Q_2$-rational divisor, lifting every closed cusp with
its full Frobenius orbit. The classes $[D-10c_0]$ and
$[C-10c_0]$ are torsion by rank zero over $\mathbb Q$ and
Manin--Drinfeld, respectively. Thus $[D-C]\in J(\mathbb Q_2)$
is torsion and reduces to zero. The verified rank condition
forces $D=C$, a contradiction.
\end{proof}
